% ─────────────────────────────────────────────────────────────────────────────
\section{Hyperparameter Choices and Ablation Factors}
\label{sec:hyperparameters}

% TODO[Kiệt]:
% Goal:
% - Explain why default hyperparameters and ablation factors were chosen.
% Flow:
% - Default configuration first, then ablation dimensions and what each tests.
% Include:
% - Keep values synchronized with Chapter 4 `tab:hparams` and Chapter 6 ablation tables.
% Avoid:
% - Reporting ablation outcomes in this setup chapter.
% Target length: 4-6 paragraphs plus existing tables.
% ─────────────────────────────────────────────────────────────────────────────

\begin{table}[H]
\centering
\caption{Default hyperparameters for AttackDRO++ (Ours).}
\label{tab:main_hyperparams}
\footnotesize
\setlength{\tabcolsep}{3.5pt}
\renewcommand{\arraystretch}{1.05}
\begin{tabularx}{\textwidth}{@{}
>{\raggedright\arraybackslash}p{0.25\textwidth}
>{\centering\arraybackslash}p{0.14\textwidth}
>{\raggedright\arraybackslash}p{0.22\textwidth}
>{\raggedright\arraybackslash}X
@{}}
\toprule
Hyperparameter & Symbol & Default value & Purpose \\
\midrule

Source attacks
& $\mathcal{A}_{\mathrm{src}}$
& PGD-20$_{\mathrm{CE}}$, DDN-$\ell_2$
& Training attack domains. \\

Number of clusters
& $K$
& 4
& Latent group granularity. \\

DRO anchor strength
& $\lambda_{\mathrm{DRO}}$
& 0.35
& Balances ERM and Cluster-DRO. \\

$q$ learning rate
& $\eta_q$
& $3 \times 10^{-4}$
& Step size for $q$ update. \\

$q$ floor
& $q_{\min}$
& 0.03
& Prevents ignored clusters. \\

$q$ warmup
& $T_w$
& 3 epochs
& Delays early reweighting. \\

Cluster refresh interval
& $R$
& 2 epochs
& K-means refresh frequency. \\

Feature bank size
& $|\mathcal{F}|$
& 4096
& Stores detached features. \\

Gradient projection dim.
& $d_{\mathrm{proj}}$
& 128
& Compresses GradFP features. \\

Gradient feature weight
& $\lambda_{\mathrm{grad}}$
& 0.1
& Weights GradFP features. \\

Label feature weight
& $\lambda_{\mathrm{lbl}}$
& 0.1
& Weights label features. \\

Clusterer
& --
& K-means
& Assigns latent groups. \\

$q$ refresh policy
& --
& Reset after refresh
& Avoids stale cluster weights. \\

\bottomrule
\end{tabularx}
\end{table}

The number of clusters is set to $K=4$ as a moderate grouping resolution. Too few
clusters may fail to capture meaningful difficulty structure, while too many clusters may
fragment the minibatch and produce noisy group-loss estimates. The default anchor
strength $\lambda_{\mathrm{DRO}}=0.35$ is chosen to keep theMulti-ATtack loss as
the dominant stabilizing signal while still allowing cluster-level adaptive correction.

The $q$-floor and warmup settings are included for stability. The floor projection
prevents any cluster from being completely ignored, while the warmup period prevents the
exponentiated-gradient update from acting on unreliable early cluster assignments. The
feature bank and refresh interval control how quickly the clusterer adapts to evolving
model representations. In the default setting, refreshing every two epochs provides a
balance between stability and adaptivity.

The gradient fingerprint is projected to $d_{\mathrm{proj}}=128$ dimensions for
computational efficiency. For CIFAR-10 with ResNet-18, the classifier-head gradient proxy
has raw dimension $C d_h = 10 \times 512 = 5120$. The fixed random projection reduces
this dimension while retaining useful information about the optimization direction
induced by each adversarial example.

Table~\ref{tab:ablation_hyperparams} lists the main hyperparameters varied in ablation
experiments. These ablations are used in Chapter~6 to justify the default configuration
and to separate the effect of the proposed components from the strong uniform
multi-attack baseline.

\begin{table}[H]
\centering
\caption{Hyperparameters varied in ablation studies.}
\label{tab:ablation_hyperparams}
\footnotesize
\setlength{\tabcolsep}{4pt}
\renewcommand{\arraystretch}{1.08}
\begin{tabularx}{\textwidth}{@{}
>{\raggedright\arraybackslash}p{0.25\textwidth}
>{\raggedright\arraybackslash}p{0.32\textwidth}
>{\raggedright\arraybackslash}X
@{}}
\toprule
Ablation factor & Values / variants & Purpose \\
\midrule

Anchor strength
& $\lambda_{\mathrm{DRO}} \in \{0.20, 0.35, 0.50\}$
& Tests the uniform-anchor trade-off. \\

Cluster count
& $K \in \{2,4,6,10\}$
& Tests sensitivity to group granularity. \\

$q$ update
& update $q$  vs. uniform $q$
& Tests adaptive reweighting. \\

Refresh schedule
& 1 epoch vs. 2 epochs
& Tests cluster-refresh frequency. \\

Gradient fingerprint
& With vs. without GradFP
& Tests optimization-level features. \\

Random clusters
& Learned vs. random assignment
& Negative control for clustering. \\

\bottomrule
\end{tabularx}
\end{table}
